\subsection{IV, 18.3--18.4 : algèbres étales et structures d’algèbre}
\label{note:iv4-unramified}
Dans 18.3.1, $B$ est déjà une $A$-algèbre de présentation finie.
Si $R=B\otimes_A B$ et $J=\ker(R\to B)$, le caractère non ramifié
équivaut à ce que la diagonale soit ouverte et fermée, donc à ce que
$J$ soit un facteur direct de $R$. Cela équivaut à la projectivité
de $B=R/J$ comme $R$-module. Aucune finitude de $B$ comme $A$-module
ne fait partie de cette équivalence. Par exemple $\mathbf Z[1/2]$
est une $\mathbf Z$-algèbre étale de présentation finie, mais n'est
pas un $\mathbf Z$-module de type fini.

Dans 18.3.2, la compatibilité $x_{in}\mapsto x_{i,n-1}$ concerne
$n\geq1$ ; la famille d'indice zéro est le choix initial.
La structure d'algèbre algébrisée en 18.3.4 comprend à la fois
l'unité $\mathcal O_X\to\mathcal B$ et la multiplication
$\mathcal B\otimes\mathcal B\to\mathcal B$. Les diagrammes d'unité,
d'associativité et de commutativité sont les identités nécessaires
pour obtenir une algèbre commutative unitaire.

Le cas local de 18.4.6 est celui d'une algèbre formellement
\emph{non ramifiée}. Dans le cas étale, $B'=A[T]/FA[T]$ est un
anneau ; le schéma correspondant est $\operatorname{Spec}(B')$.
La partie de 18.4.6 utilisée deux fois dans 18.4.7 est (i), qui fournit
la surjection locale. Dans 18.4.8, les applications du recouvrement
sont $g'_\alpha$.

Pour le changement de base de 18.4.14, l'identité des différentielles est
\[
 \Omega^1_{(B\otimes_A k)/k}
   \simeq \Omega^1_{B/A}\otimes_B(B\otimes_A k)
   \simeq \Omega^1_{B/A}/\mathfrak m\Omega^1_{B/A}.
\]
Le corps $k=A/\mathfrak m$ n'a pas, à lui seul, de structure canonique
de $B$-module. C'est le module sur $B/\mathfrak mB$ qui intervient.
La finitude de $\Omega^1_{B/A}$ et l'inclusion de $\mathfrak mB$ dans
le radical de $B$ permettent alors l'application de Nakayama.

\subsection{IV, 18.10--18.12 : base et parenthèses des arguments}
\label{note:iv4-normalization}
En 18.10.14, après extension de $K$ à $K'$, la fermeture intégrale
est prise sur $Y'$ dans $L\otimes_K K'$. La décomposition de cette
algèbre en facteurs donne notamment la fermeture de $Y'$ dans
$L_1$ ; le résultat applicable est 18.10.13.
En 18.11.2, le module de type fini est $\Omega^1_{C/k}$,
identifié à son module complété ; 18.11.1 est la référence qui
fournit la finitude. La référence à 0, 20.6.21(i) appartient à la
parenthèse sur la suite exacte des trois corps.
Dans 18.12.15, la comparaison avec 6.14.4, partie II, se rapporte
au passage à une sous-algèbre de type fini.

\subsection{IV, 19.1--19.3 et 19.7--19.9 : immersions et fibres}
\label{note:iv4-immersions}
Le module conormal $J/J^2$ est naturellement un module sur $A/J$,
puisque $J$ l'annule. C'est sur cet anneau qu'il est libre en 19.1.4,
de rang égal à la codimension transversale. Ses puissances graduées
ont alors les rangs binomiaux utilisés dans le calcul de Hilbert--Samuel.
En 19.1.7, les propriétés des immersions sont celles de 19.1.5(i)--(iv) ;
la seconde assertion visée est celle de (ii).

Le critère de régularité relative utilisé à la fin de 19.2.4 est
0, 15.1.16. Pour une immersion ouverte dans 19.2.7(i), l'idéal est nul,
mais la régularité transversale relative exige encore la platitude
de la source sur $S$. L'identité d'un $S$-schéma non plat est déjà
un exemple qui montre la nécessité de cette condition.
En 19.3.4, les immersions sont $X\to Y$ et $X'\to Y'$ :
le spectre de l'algèbre quotient est le sous-schéma du spectre de
l'algèbre de polynômes.

En 19.7.5(ii), l'hypothèse sur $X$ fournit la régularité de la suite
sur $\operatorname{gr}_{\mathcal K_z}(\mathcal O_{X,z})$.
Elle ne fournit pas directement cette régularité sur le gradué
d'un module arbitraire $\mathcal F_z$ ; c'est ensuite le critère
19.7.1 qui porte sur ce module. Dans 19.9.8, les deux cas du
raisonnement utilisent respectivement les bornes de profondeur
$1$ et $2$ dans la définition de $Z'_\lambda$.

\subsection{IV, 19.4.6 : platitude et changement de base d’un éclatement}
\label{note:iv4-blowups}
Travaillons sur $S=\operatorname{Spec}(A)$, $X=\operatorname{Spec}(B)$,
avec un idéal $J\subset B$. Posons
$Q_m=B/J^m$ pour $m\geq1$ et $G_r=J^r/J^{r+1}$ pour $r\geq0$.
La platitude de $\bigoplus_rG_r$ équivaut à celle de tous ses termes.
Puisque $Q_1=G_0$, les suites
\[
 0\longrightarrow G_r\longrightarrow Q_{r+1}
   \longrightarrow Q_r\longrightarrow0
\]
montrent que cette condition équivaut à la platitude de tous les $Q_m$.
On utilise dans un sens la stabilité de la platitude par extension,
et dans l'autre la platitude du noyau d'une surjection entre modules plats.

Si $Q_m$ est plat, la suite des Tor associée à
$0\to J^m\to B\to Q_m\to0$ donne l'injectivité de
$J^m\otimes_A A'\to B\otimes_A A'$ pour toute $A$-algèbre $A'$.
Pour la réciproque, supposons en plus $B$ plat sur $A$.
La même suite donne alors $\operatorname{Tor}_1^A(Q_m,A')=0$.
Prendre l'algèbre à idéal de carré nul $A'=A\oplus M$, pour un
$A$-module arbitraire $M$, donne
$\operatorname{Tor}_1^A(Q_m,M)=0$, donc la platitude de $Q_m$.
Sans la platitude de $B$, cette réciproque est fausse : pour $J=0$,
les applications sur les puissances positives de $J$ sont injectives
quel que soit $B$, y compris lorsque $B$ n'est pas plat.

Les conclusions sur l'éclatement et son diviseur exceptionnel se
prennent sous la platitude des $Q_m$, ou de manière équivalente du
gradué. L'injectivité précédente identifie chaque $J^m\otimes_A A'$
à la puissance de l'idéal image dans $B\otimes_A A'$.
L'algèbre de Rees commute donc à ce changement de base ; la propriété
correspondante de $\operatorname{Proj}$ donne l'isomorphisme des
éclatements annoncé dans l'énoncé.

\subsection{IV, 19.7.1 : relations de Koszul et filtration}
\label{note:iv4-koszul}
Posons $Q=\mathfrak R^a/\mathfrak R^{a+1}$ et $I=(f_1,\ldots,f_n)$.
La suite $\mathbf f$ est $Q$-régulière. Une relation
$\sum_i a_if_i=0$ dans $Q$ ne force pas la nullité de chaque $a_i$ :
la relation $(f_2,-f_1)$ pour deux éléments est l'exemple élémentaire.
Il faut modifier les coefficients par une relation de Koszul.

Voici le contrôle de l'ordre des coefficients. Soit
$u\in I^sK_1(\mathbf f;Q)$ un cycle, où $s=b-1\geq1$.
L'exactitude du complexe de Koszul donne $u=\delta_2C$.
La quasi-régularité donne
$\operatorname{gr}_I(Q)\simeq(Q/IQ)[T_1,\ldots,T_n]$.
Le complexe de Koszul des indéterminées est exact en degrés positifs
sur ce module polynomial.

Si $C\in I^dK_2$ et $d<s-1$, le terme initial de $C$ est un cycle,
car $\delta_2C\in I^sK_1$ et la différentielle augmente de un le
degré des coefficients. Pour $d=0$, ce cycle est nul : une préimage
dans $K_3$ devrait avoir un degré de coefficient négatif. Pour
$d\geq1$, il est le bord d'une classe de degré $d-1$ ; en relevant
cette classe en $D\in I^{d-1}K_3$, on remplace $C$ par
$C-\delta_3D\in I^{d+1}K_2$, sans changer $\delta_2C$.
Cette procédure finie donne $C\in I^{s-1}K_2$.

Les coefficients de $C$ se relèvent en
$c_{ij}\in\mathfrak R^a I^{b-2}$ avec
$c_{ij}=-c_{ji}$ et $c_{ii}=0$. Remplacer $a_i$ par
$a_i-\sum_jc_{ij}f_j$ conserve $\sum_i a_if_i$, puisque les termes
antisymétriques s'annulent. Les nouveaux coefficients appartiennent
à $\mathfrak R^a I^{b-1}\cap\mathfrak R^{a+1}$, comme requis
pour la récurrence de 19.7.1.16. Pour $b=1$, la relation considérée
est immédiate. Aucun passage à une filtration complétée n'intervient.

\subsection{IV, 21.10.16 : longueurs après changement de base}
\label{note:iv4-lengths}
Avec les anneaux locaux de la démonstration, posons
$A'_1=A'\otimes_A A_1$ et $M=A'_1/\mathfrak m_{x'}A'_1$.
La décomposition d'un module de longueur finie sur l'algèbre
semi-locale finie $A'_1$ donne
\[
 \operatorname{long}_{A_1}(M)
 =\sum_{x'_1}\operatorname{long}_{(A'_1)_{x'_1}}(M_{x'_1})
       [k(x'_1):k(x_1)].
\]
Chaque facteur simple au point $x'_1$ est le corps résiduel $k(x'_1)$,
dont la longueur sur $A_1$ est son degré sur $k(x_1)$.
Le membre de droite est donc le premier membre de 21.10.16.2.

Comme $A_1$ est plat sur $A$ et
$M=(A'/\mathfrak m_{x'}A')\otimes_A A_1$, une suite de composition
sur $A$ reste exacte après tensorisation. Elle donne
\[
 \operatorname{long}_{A_1}(M)
 =[k(x'):k(x)]\,
   \operatorname{long}_{A_1}(A_1/\mathfrak m_xA_1).
\]
C'est l'égalité voulue. Elle ne requiert pas que l'extension
$k(x_1)/k(x)$ soit finie et n'utilise aucun quotient de degrés
qui pourrait ne pas être défini.
\par
